\documentclass[11pt]{article}

\usepackage[margin=1in]{geometry}
\usepackage[T1]{fontenc}
\usepackage{lmodern}
\usepackage{microtype}
\usepackage{hyperref}
\usepackage{enumitem}
\usepackage{amsmath, amssymb}
\usepackage{booktabs}
\usepackage{xcolor}
\usepackage{titlesec}

\hypersetup{
  colorlinks=true,
  linkcolor=blue!45!black,
  urlcolor=blue!45!black,
  citecolor=blue!45!black
}

\setlist[itemize]{leftmargin=*, itemsep=2pt, topsep=4pt}
\setlist[enumerate]{leftmargin=*, itemsep=2pt, topsep=4pt}

\titleformat{\section}{\Large\bfseries}{\thesection}{0.6em}{}
\titleformat{\subsection}{\large\bfseries}{\thesubsection}{0.6em}{}
\titleformat{\subsubsection}{\normalsize\bfseries}{\thesubsubsection}{0.6em}{}

\title{Expansion Copilot (Sydney Gyms)\\\large Demo Build Report (Pre-Transport)}
\author{}
\date{}

\begin{document}
\maketitle

\noindent\textbf{Tagline}: Turn place data into smarter expansion decisions.\\
\textbf{Core question}: \emph{Where should I open my next gym in Sydney, and why?}\\
\textbf{Scope}: Sydney, Australia \textbullet\ Boutique / mid-market gyms \textbullet\ Decision support (not prediction)

\bigskip
\hrule
\bigskip

\section{What we built (so far)}

\subsection{Product}
\begin{itemize}
  \item A Streamlit MVP that:
  \begin{itemize}
    \item shows \textbf{ranked candidate areas}
    \item visualizes them on a map (color = score)
    \item provides a \textbf{``Why here?''} explanation panel with transparent drivers
    \item supports \textbf{Compare} for 2--3 areas
    \item includes \textbf{guardrails} and \textbf{judge-mode presets} for demo stability
  \end{itemize}
\end{itemize}

\subsection{Data + feature pipeline (offline, reproducible)}
\begin{itemize}
  \item Reads Foursquare Open Source Places (Parquet) subset for Sydney.
  \item Builds an explainable feature table aggregated by area (H3).
  \item Computes an Opportunity Score with a transparent weighted formula.
  \item Writes \textbf{processed artifacts} consumed by the app (no recompute in the UI).
\end{itemize}

\section{Repository structure (relevant parts)}
\begin{itemize}
  \item \texttt{app/streamlit\_app.py}: current demo app (single-file Streamlit app, upgraded with demo controls)
  \item \texttt{app/app.py} + \texttt{src/app/*}: modular Streamlit MVP (cleaner structure, same data contract)
  \item \texttt{src/features/build\_opportunity\_features.py}: feature engineering for area-level metrics
  \item \texttt{src/scoring/opportunity\_score.py}: scoring math (normalization, weights, breakdown)
  \item \texttt{scripts/run\_build\_features.py}: produces processed outputs
  \item \texttt{data/interim/}: Sydney subset + cleaned categories used to build features
  \item \texttt{data/processed/}: feature table + ranked areas (app reads these)
\end{itemize}

\section{Data inputs (pre-transport)}

\subsection{Foursquare Open Source Places (Sydney subset)}
\textbf{Inputs (from prior tasks):}
\begin{itemize}
  \item \texttt{data/interim/places\_sydney.parquet}
  \item \texttt{data/interim/categories\_clean.parquet}
\end{itemize}

\noindent We explicitly designed the pipeline to \textbf{not assume optional columns exist}. Where columns are missing (e.g., refresh timestamps), we use safe fallbacks or leave placeholders.

\section{Core feature engineering (area-level)}
\noindent We aggregate to an \textbf{area unit} (v1: H3 at resolution 8).

\subsection{Place $\rightarrow$ role flags (competitor / complement / excluded)}
We classify each place into business roles using a transparent gym taxonomy:
\begin{itemize}
  \item competitors: gyms / yoga / pilates / martial arts (configurable)
  \item complements: healthy cafes, physio, wellness, parks, sports goods, etc.
  \item exclusions: categories that should not count toward active commercial ecosystem
\end{itemize}

\noindent\textbf{Implementation}: \texttt{src/features/category\_mapping.py}\\
\textbf{Key design constraint}: explode/join category IDs \textbf{without inflating places}.

\subsubsection*{Method}
\begin{enumerate}
  \item explode \texttt{category\_ids} to one row per (\texttt{place\_id}, \texttt{category\_id})
  \item join to cleaned categories table for \texttt{category\_name} and \texttt{category\_label}
  \item keyword-match into flags
  \item group back to \texttt{place\_id} using \texttt{any()} for boolean flags
\end{enumerate}

\noindent\textbf{Output per place:}
\begin{itemize}
  \item \texttt{is\_competitor}, \texttt{is\_complementary}, \texttt{is\_commercial\_other}, \texttt{is\_excluded}
  \item \texttt{unique\_category\_count} (per place, later summed by area)
\end{itemize}

\subsection{Derived POI Data Quality Score (0..1)}
\noindent\textbf{Implementation}: \texttt{src/features/data\_quality.py}

\bigskip
\noindent We compute a transparent POI score, clamped to $[0, 1]$, based on field presence.

\paragraph{Base score}
\[
q = 1.0
\]

\paragraph{Penalties (applied if columns exist)}
\begin{itemize}
  \item missing coords: $0.40$
  \item missing name: $0.20$
  \item missing category: $0.15$
  \item unresolved flags present: $0.25$
\end{itemize}

\paragraph{Formula}
Let indicator functions evaluate to 1 when true, else 0:
\[
q = 1
 - 0.40\cdot \mathbb{1}[\text{lat or lon missing}]
 - 0.20\cdot \mathbb{1}[\text{name missing}]
 - 0.15\cdot \mathbb{1}[\text{category missing}]
 - 0.25\cdot \mathbb{1}[\text{unresolved\_flags present}]
\]

\paragraph{Clamp}
\[
q := \min(1, \max(0, q))
\]

\paragraph{Area-level data quality}
\[
\text{mean\_data\_quality\_score} = \frac{1}{N}\sum_{i=1}^{N} q_i
\]

\subsection{Active POIs}
In v1, we treat ``active'' as:
\begin{itemize}
  \item not excluded by taxonomy (\texttt{\textasciitilde is\_excluded})
  \item (optionally) not closed if \texttt{date\_closed} exists (not currently in our canonical inputs)
\end{itemize}

\noindent So:
\[
\text{is\_active} = \neg \text{is\_excluded}
\]

\subsection{Aggregates per area (counts)}
\noindent\textbf{Implementation}: \texttt{src/features/build\_opportunity\_features.py}

\noindent For each \texttt{area\_id} we compute:
\begin{itemize}
  \item \texttt{total\_poi\_count}
  \item \texttt{active\_poi\_count}
  \item \texttt{competitor\_count}
  \item \texttt{complementary\_count}
  \item \texttt{commercial\_other\_count}
  \item \texttt{unique\_category\_count} (summed)
  \item \texttt{mean\_data\_quality\_score} (mean)
\end{itemize}

\subsection{Derived area features (transparent formulas)}
We compute the core proxies (with $\varepsilon = 1.0$ for numerical stability).

\paragraph{Competitor density proxy}
\[
\text{competitor\_density\_proxy}
=
\frac{\text{competitor\_count}}{\text{active\_poi\_count} + \varepsilon}
\]

\paragraph{Complementarity ratio}
\[
\text{complementarity\_ratio}
=
\frac{\text{complementary\_count}}{\text{active\_poi\_count} + \varepsilon}
\]

\paragraph{Commercial activity proxy}
\[
\text{commercial\_activity\_proxy} = \text{active\_poi\_count}
\]

\paragraph{Saturation proxy}
\[
\text{saturation\_proxy}
=
\text{competitor\_density\_proxy}\cdot (1 - \text{complementarity\_ratio})
\]

\noindent\texttt{recent\_refresh\_share} is currently a placeholder because it depends on refresh timestamp columns that may not exist in the dataset we’re using.

\section{Opportunity Score (math + implementation)}
\noindent\textbf{Implementation}: \texttt{src/scoring/opportunity\_score.py}

\subsection{Stabilization (log1p)}
For heavy-tailed signals we apply log transforms:
\begin{itemize}
  \item $\texttt{commercial\_activity\_proxy\_log}=\log(1+\texttt{commercial\_activity\_proxy})$
  \item $\texttt{diversity\_log}=\log(1+\texttt{unique\_category\_count})$
  \item $\texttt{competition\_log}=\log(1+\texttt{competitor\_density\_proxy})$
  \item $\texttt{saturation\_log}=\log(1+\texttt{saturation\_proxy})$
\end{itemize}

\subsection{Normalization (safe min--max)}
For each signal $x$:
\[
n(x)=
\begin{cases}
0 & \text{if } \max(x)-\min(x) \le 10^{-9}\\
\frac{x-\min(x)}{\max(x)-\min(x)} & \text{otherwise}
\end{cases}
\]

\subsection{Weighted sum (transparent)}
Default weights (\texttt{OpportunityWeights}):
\begin{itemize}
  \item complementarity: $+1.2$
  \item commercial activity: $+0.8$
  \item diversity: $+0.6$
  \item direct competition: $-1.0$
  \item saturation: $-0.7$
  \item data quality: $+0.5$
\end{itemize}

\noindent Let normalized features be:
\begin{itemize}
  \item $n_c$ = \texttt{n\_complementarity}
  \item $n_a$ = \texttt{n\_activity}
  \item $n_d$ = \texttt{n\_diversity}
  \item $n_k$ = \texttt{n\_competition}
  \item $n_s$ = \texttt{n\_saturation}
  \item $n_q$ = \texttt{n\_quality}
\end{itemize}

\noindent Then:
\[
\text{OpportunityScore}
=
1.2n_c
+0.8n_a
+0.6n_d
-1.0n_k
-0.7n_s
+0.5n_q
\]

\subsection{Reasons (built-in explanation function)}
We generate \texttt{reasons\_top3} via \texttt{top\_reasons(row)}:
\begin{itemize}
  \item compute contribution components (e.g., \texttt{score\_complementarity}, \texttt{score\_activity}, etc.)
  \item return the top 3 labels by magnitude
\end{itemize}
This is deterministic and demo-safe (no LLM dependency).

\section{Guardrails (demo stability + honesty)}
\noindent In \texttt{scripts/run\_build\_features.py} we apply:
\begin{itemize}
  \item \texttt{active\_poi\_count >= 200} to prevent tiny cells dominating normalization.
\end{itemize}

\noindent In the Streamlit app we also expose:
\begin{itemize}
  \item minimum active POIs
  \item minimum mean data quality
\end{itemize}

\noindent And we label outputs as \textbf{decision support}, not revenue/demand forecasting.

\section{Processed outputs (what the app reads)}
\noindent Produced by \texttt{scripts/run\_build\_features.py}:
\begin{itemize}
  \item \texttt{data/processed/sydney\_area\_features.parquet}
  \item \texttt{data/processed/sydney\_ranked\_areas.parquet}
  \item \texttt{data/exports/top\_ranked\_areas.csv}
  \item docs previews:
  \begin{itemize}
    \item \texttt{docs/feature\_summary.md}
    \item \texttt{docs/top\_areas\_preview.md}
  \end{itemize}
\end{itemize}

\noindent The Streamlit demo is designed to \textbf{prefer precomputed ranked outputs} for performance.

\section{Streamlit MVP (demo behavior)}

\subsection{Required sections implemented}
\begin{itemize}
  \item Landing/setup panel (Sydney + Gym fixed for v1)
  \item Opportunity map (tooltip shows score drivers)
  \item Top candidates panel (table + pinning)
  \item ``Why this location?'' (metrics + narrative + examples)
  \item Compare mode (2--3 areas)
  \item Demo narrative mode (limitations + how score works)
\end{itemize}

\subsection{Judge-mode enhancements}
\begin{itemize}
  \item Judge Mode toggle + hero buttons (Map / Why / Compare)
  \item Stable presets for top\_k and guardrails
  \item Clear closing line: \emph{Businesses do not need more map data. They need fewer bad location decisions.}
\end{itemize}

\section{What changes when we add transport data (not implemented yet)}
Transport data would add an \textbf{accessibility feature} (not demand):
\begin{itemize}
  \item \texttt{transit\_access\_score} per area (0..1), derived from station entrances / stops density and proximity.
\end{itemize}

\noindent It would enter the score as a new positive term:
\[
\text{OpportunityScore}_{new}
= \text{OpportunityScore} + w_t\cdot n(\text{transit\_access\_score})
\]

\noindent Where $w_t$ is a modest weight (e.g., 0.3--0.7) so it acts as a tie-breaker and improves ``Why here?'' defensibility without overpowering the ecosystem logic.

\end{document}

